\chapter{Vervollständigung, Fixpunkte und Kategorie}
\label{b47guide:ch:completion-fixed-points}

\section{Fortsetzungen metrischer Abbildungen}

\begin{theorem}[Gleichmäßig stetige Fortsetzung]
\label{b47guide:thm:uniform-extension}
Seien $X,Y$ metrische Räume, sei $D\subseteq X$ dicht, sei $Y$ vollständig und sei
$f:D\to Y$ gleichmäßig stetig. Dann gibt es genau eine stetige
Fortsetzung $F:X\to Y$ von $f$. Diese Fortsetzung ist gleichmäßig
stetig. Für $x\in X$ und jede Folge $(a_n)$ in $D$ mit $a_n\to x$ gilt
\[
 F(x)=\lim_{n\to\infty}f(a_n).
\]
Ist $f$ $L$-Lipschitz, so ist $F$ ebenfalls $L$-Lipschitz.
Ist $f$ eine isometrische Einbettung, so ist auch $F$ eine
isometrische Einbettung.
\cite{baire-notes}
\end{theorem}
\begin{theorem}[McShane--Whitney: Lipschitz-Fortsetzung auf den ganzen Raum]
\label{b47guide:thm:mcshane-whitney}
Sei $(X,d)$ ein metrischer Raum, sei $\varnothing\ne A\subseteq X$ und
sei $f:A\to\mathbb R$ $L$-Lipschitz mit $L\geq0$. Für $x\in X$ setzen wir
\[
\begin{aligned}
 F_-(x)&=\sup_{a\in A}\bigl(f(a)-L\,d(x,a)\bigr),\\
 F_+(x)&=\inf_{a\in A}\bigl(f(a)+L\,d(x,a)\bigr).
\end{aligned}
\]
Beide Formeln haben endliche reelle Werte und definieren
$L$-Lipschitz-Fortsetzungen von $f$ auf $X$.
Jede weitere $L$-Lipschitz-Fortsetzung $F:X\to\mathbb R$ erfüllt
$F_-(x)\leq F(x)\leq F_+(x)$ für alle $x\in X$.
Die Menge $A$ muss weder dicht noch abgeschlossen sein; Vollständigkeit
von $X$ wird nicht vorausgesetzt. Für $A=\varnothing$ ist jede konstante
Abbildung eine geeignete Fortsetzung, ohne Verwendung dieser Formeln.
\cite{lipschitz-extension}
\end{theorem}
\section{Der vollständige Raum zu einer Metrik}

\begin{definition}[Vervollständigung]
\label{b47guide:def:completion}
Eine \emph{Vervollständigung} eines metrischen Raums $(X,d)$ ist ein
Paar $(\widehat X,i)$ aus einem vollständigen metrischen Raum
$(\widehat X,\widehat d)$ und einer isometrischen Einbettung
$i:X\to\widehat X$, deren Bild in $\widehat X$ dicht ist.
\end{definition}

\begin{theorem}[Existenz und eindeutige Bestimmtheit der Vervollständigung]
\label{b47guide:thm:completion-existence}
Jeder metrische Raum besitzt eine Vervollständigung.
Sind $(\widehat X,i)$ und $(\widetilde X,j)$ zwei Vervollständigungen
von $X$, so gibt es genau eine Isometrie
$U:\widehat X\to\widetilde X$ mit
\[
 U\circ i=j.
\]
Die Eindeutigkeit bezieht sich auf die vorgegebenen Einbettungen.
\cite{completion-notes}
\end{theorem}

\begin{theorem}[Vervollständigung als Abschluss]
\label{b47guide:thm:completion-closure}
Ist $j:X\to Z$ eine isometrische Einbettung in einen vollständigen
metrischen Raum $Z$, so ist $(\overline{j[X]}^{\,Z},j)$ eine
Vervollständigung von $X$, wobei die Zielmenge von $j$ auf diesen
Abschluss eingeschränkt wird.
Für jede Vervollständigung $(\widehat X,i)$ gilt
\[
 X\text{ vollständig}\quad\Longleftrightarrow\quad i[X]=\widehat X.
\]
Außerdem ist $X$ genau dann total beschränkt, wenn $\widehat X$
kompakt ist.
\cite{completion-notes}
\end{theorem}

\begin{theorem}[Fortsetzungscharakterisierung]
\label{b47guide:thm:completion-universal}
Sei $(\widehat X,i)$ eine Vervollständigung von $X$ und sei $Y$
ein vollständiger metrischer Raum. Jede gleichmäßig stetige
Abbildung $f:X\to Y$ besitzt genau eine stetige Abbildung
$\widehat f:\widehat X\to Y$ mit $\widehat f\circ i=f$.
Sie ist gleichmäßig stetig; eine vorhandene Lipschitz-Konstante
von $f$ bleibt erhalten.
\cite{completion-notes}
\end{theorem}

\section{Kontraktionen und quantitative Fixpunktsätze}

\begin{definition}[Kontraktion und Fixpunktiteration]
\label{b47guide:def:contraction}
Sei $(X,d)$ ein metrischer Raum. Eine Abbildung $T:X\to X$ heißt
\emph{Kontraktion}, wenn es eine reelle Zahl $0\leq q<1$ gibt mit
\[
 d(Tx,Ty)\leq q\,d(x,y)\qquad(x,y\in X).
\]
Ein Punkt $p\in X$ mit $Tp=p$ heißt \emph{Fixpunkt}.
Zu einem Startpunkt $x_0\in X$ heißt die durch $x_{n+1}=Tx_n$
definierte Folge \emph{Fixpunktiteration}.
\end{definition}

\begin{theorem}[Banachscher Fixpunktsatz]
\label{b47guide:thm:banach}
Jede Kontraktion eines \emph{nichtleeren vollständigen} metrischen
Raums in sich besitzt genau einen Fixpunkt $p$.
Für jeden Startpunkt $x_0$ konvergiert die Fixpunktiteration gegen $p$.
\cite{banach-notes}
\end{theorem}

\begin{theorem}[Fehlerabschätzungen der Fixpunktiteration]
\label{b47guide:thm:banach-errors}
Unter den Voraussetzungen des Banachschen Fixpunktsatzes gelten
für $n\geq1$ die Abschätzungen
\[
\begin{aligned}
 d(x_n,p)&\leq \frac{q^n}{1-q}\,d(x_1,x_0)
 &&\text{(a priori)},\\
 d(x_n,p)&\leq \frac{q}{1-q}\,d(x_n,x_{n-1})
 &&\text{(a posteriori)}.
\end{aligned}
\]
Für jeden Punkt $x\in X$ gilt zusätzlich
\[
 d(x,p)\leq\frac{d(x,Tx)}{1-q}.
\]
Diese Formeln gelten auch für $q=0$; dann ist $x_1=p$.
\cite{banach-notes}
\end{theorem}

\begin{theorem}[Eine kontrahierende Iterierte genügt]
\label{b47guide:thm:iterated-contraction}
Sei $X$ ein nichtleerer vollständiger metrischer Raum und sei
$T:X\to X$ eine Abbildung. Ist für ein $m\geq1$ die $m$-fache
Iteration $T^m$ eine Kontraktion, so besitzt $T$ genau einen
Fixpunkt $p$. Für jedes $x\in X$ gilt $T^n x\to p$.
Eine gesonderte Stetigkeitsannahme an $T$ ist nicht nötig.
\end{theorem}

\begin{theorem}[Stabilität von Fixpunkten]
\label{b47guide:thm:fixed-point-stability}
Sei $X$ nichtleer und vollständig. Seien $T,S:X\to X$
Kontraktionen mit Fixpunkten $p_T,p_S$, und sei $q<1$
eine nichtnegative Lipschitz-Konstante von $T$.
Gilt für eine reelle Zahl $\eta\geq0$
\[
 d(Tx,Sx)\leq\eta\qquad(x\in X),
\]
so gilt
\[
 d(p_T,p_S)\leq\frac{\eta}{1-q}.
\]
\end{theorem}

\begin{theorem}[Edelsteins Fixpunktsatz auf kompakten metrischen Räumen]
\label{b47guide:thm:edelstein}
Sei $(X,d)$ ein nichtleerer kompakter metrischer Raum und sei $T:X\to X$
eine Abbildung mit
\[
 d(Tx,Ty)<d(x,y)\qquad(x,y\in X,\ x\ne y).
\]
Dann besitzt $T$ genau einen Fixpunkt $p$. Für jeden Startpunkt $x\in X$
gilt $T^n x\to p$.
Ein gemeinsamer Kontraktionsfaktor $q<1$ und eine gesonderte
Stetigkeitsannahme werden nicht verlangt. Die geometrischen
Fehlerabschätzungen des Banachschen Satzes folgen aus diesen
Voraussetzungen allein nicht. \cite{edelstein}
\end{theorem}
\section{Der Bairesche Kategoriensatz}

\begin{definition}[Nirgends dichte und magere Mengen]
\label{b47guide:def:category}
Eine Teilmenge $A$ eines metrischen Raums $X$ heißt
\emph{nirgends dicht in $X$}, wenn der Abschluss $\overline A^{\,X}$
keine nichtleere offene Teilmenge von $X$ enthält.
Eine Menge $M\subseteq X$ heißt \emph{mager in $X$}, wenn sie
Vereinigung abzählbar vieler in $X$ nirgends dichter Mengen ist.
Alle diese Begriffe beziehen sich auf den angegebenen Umgebungsraum.
\end{definition}

\begin{theorem}[Bairescher Kategoriensatz für vollständige metrische Räume]
\label{b47guide:thm:baire}
Sei $X$ ein vollständiger metrischer Raum und seien $G_n\subseteq X$
für $n\geq1$ offen und dicht. Dann ist
\[
 \bigcap_{n\geq1}G_n
\]
dicht in $X$. Äquivalent: Keine nichtleere offene Teilmenge von $X$
ist mager in $X$.
\cite{baire-notes}
\end{theorem}

\begin{theorem}[Folgerungen des Kategoriensatzes]
\label{b47guide:thm:baire-consequences}
Ein nichtleerer vollständiger metrischer Raum ist nicht mager in sich.
Wird er als abzählbare Vereinigung abgeschlossener Mengen geschrieben,
so besitzt wenigstens eine dieser Mengen nichtleeres Inneres.
Ein nichtleerer vollständiger metrischer Raum ohne isolierte Punkte
ist überabzählbar; dabei heißt $x$ \emph{isoliert}, wenn $\{x\}$ offen ist.
\end{theorem}